\begin{exercise}[Enthalpie libre d'un gaz parfait]
\begin{questions}
\question Rappeler l'expression de l'enthalpie libre G(T,P) de $n ~ \si{\mole}$ de gaz parfait, maintenues à $T$ constante.

\question Indiquer comment la relation de Gibbs et Helmholtz permet de montrer que l’enthalpie $H$ du gaz parfait ne dépend que de la température $T$.

\question  En déduire que $U$ ne dépend que de $T$.

\question Rappeler les deux lois de Joule en reliant les différentielles $dU$ et $dH$ à la variation $dT$.
\end{questions}

\end{exercise}